\scanpage{96}
\mspage{24}
8. Une façon d'obtenir une configuration ``polygône régulier à $n$ côtés'' \inserted{$P_n$} dans une sphère complexe \inserted{$\Sigma$}, \deleted{est} i.e. dans une droite projective complexe, est de partir d'un polygone régulier à $2n$ côtés $\widetilde P_{2n}$ (\unsure{moralement}, un revêtement de degré $2$ de $P_n$), plongé dans un \deleted{espace} \inserted{plan} vectoriel réel $V_{\widetilde P_{2n}}$ (en tant que polygône régulier) et de poser

\margin{$P_n=P=(S,A,R\subset S\times A)$}

\[
\tag{1}\Sigma=\mathbb P(V_{\widetilde P_{2n}})(\mathbb C)
\]
en prenant
\[
\tag{2}P_n=\widetilde P_{2n}/\{\pm1\}
\]
\margin{\uline{NB} $-1$ opère sur le polygone combinatoire $P_{2n}$ par antipodisme}

\deleted{et en} donc
\[
S=\widetilde S/\pm1,\quad A=\widetilde A/\pm1,\quad R=\widetilde R/\pm1,
\]
et en définissant le plongement
\[
\tag{3}S\hookrightarrow\Sigma
\]
par passage au quotient en termes de
(
\diagram{\begin{tikzcd}[ampersand replacement=\&]
\widetilde S\arrow[r]\&\Sigma\\
R_{\tilde s}\arrow[r,mapsto]\&
\mathbb C.R_{\tilde s}\in\mathbb P(V_{\mathbb C})\unsure{\{=|(=}\mathbb P(V)(\mathbb C).
\end{tikzcd}}

\deleted{De cette} En fait, \deleted{L'application (3) se \unsure{factorise}} \inserted{par construction $\Sigma$ provient} de la droite projective \uline{réelle} $\mathbb P(\underset{V}{\underline{V_{\widetilde P_{2n}}}})$, et l'application (3) se factorise par \deleted{$\Sigma_{\mathbb R}\Vert$}
\[
\tag{4}\Sigma_{\mathbb R}=\boldsymbol\Sigma_{\widetilde P_{2n}}(\mathbb R)
\]
où on a posé
\[
\tag{5}
\underset{\substack{\Vert\\\boldsymbol\Sigma_{P_n}}}{\boldsymbol\Sigma_{\widetilde P_{2n}}}
=\mathbb P(V_{\widetilde P_{2n}})
\quad\text{(ne dépend en fait que de }P_n=\widetilde P_{2n}/\pm1\text{)}
\]

\margin{95}
